\documentclass{article}
\usepackage{amsmath,amssymb,amsthm,algorithm,algorithmic}
\usepackage{bm}
\newtheorem{theorem}{Theorem}
\newtheorem{lemma}{Lemma}
\newtheorem{assumption}{Assumption}
\begin{document}

\section*{Primal–Dual Hybrid Gradient (PDHG) Algorithm}

\section*{Mathematical Formulation}
Consider the convex–concave saddle‑point problem  

\[
\min_{x\in\mathbb{R}^n}\max_{y\in\mathbb{R}^m}
\;\mathcal{L}(x,y)\;:=\;f(x)+\langle Kx,y\rangle-g^{*}(y),
\tag{1}
\]

where  

\begin{itemize}
  \item $f:\mathbb{R}^{n}\!\to\!\mathbb{R}\cup\{+\infty\}$ is proper, closed and convex (possibly non‑smooth);
  \item $g^{*}:\mathbb{R}^{m}\!\to\!\mathbb{R}\cup\{+\infty\}$ is the convex conjugate of a proper, closed, convex $g$;
  \item $K\in\mathbb{R}^{m\times n}$ is a linear operator.
\end{itemize}

Define the proximal operators  

\[
\operatorname{prox}_{\tau f}(v):=\arg\min_{x}\left\{f(x)+\frac{1}{2\tau}\|x-v\|^{2}\right\},
\qquad
\operatorname{prox}_{\sigma g^{*}}(u):=\arg\min_{y}\left\{g^{*}(y)+\frac{1}{2\sigma}\|y-u\|^{2}\right\}.
\]

Given stepsizes $\tau>0$, $\sigma>0$ satisfying  

\[
\tau\sigma\|K\|^{2}<1,
\tag{2}
\]

the PDHG updates are  

\[
\boxed{
\begin{aligned}
y_{k+1}&=\operatorname{prox}_{\sigma g^{*}}\!\bigl(y_{k}+\sigma K\bar{x}_{k}\bigr),\\[2mm]
x_{k+1}&=\operatorname{prox}_{\tau f}\!\bigl(x_{k}-\tau K^{\!\top}y_{k+1}\bigr),\\[2mm]
\bar{x}_{k+1}&=x_{k+1}+\theta\,(x_{k+1}-x_{k}),
\end{aligned}}
\tag{3}
\]

with extrapolation parameter $\theta\in[0,1]$ (standard choice $\theta=1$).  
The algorithm keeps the **dual variable** $y_{k}$ and the **primal variable** $x_{k}$; $\bar{x}_{k}$ is an extrapolated primal point used only in the dual update.

\section*{Pseudocode}
\begin{algorithm}[H]
\caption{Primal–Dual Hybrid Gradient (PDHG)}\label{alg:PDHG}
\begin{algorithmic}[1]
\REQUIRE $f$, $g^{*}$, linear operator $K$, stepsizes $\tau,\sigma>0$, extrapolation $\theta\in[0,1]$, initial $x_{0},y_{0}$.
\STATE $\bar{x}_{0}\gets x_{0}$
\FOR{$k=0,1,2,\dots$}
    \STATE $y_{k+1}\gets\operatorname{prox}_{\sigma g^{*}}\!\bigl(y_{k}+\sigma K\bar{x}_{k}\bigr)$ \hfill // dual step
    \STATE $x_{k+1}\gets\operatorname{prox}_{\tau f}\!\bigl(x_{k}-\tau K^{\!\top}y_{k+1}\bigr)$ \hfill // primal step
    \STATE $\bar{x}_{k+1}\gets x_{k+1}+\theta\,(x_{k+1}-x_{k})$ \hfill // extrapolation
\ENDFOR
\ENSURE Approximate saddle point $(x_{k},y_{k})$.
\end{algorithmic}
\end{algorithm}

\section*{Dry Run Example}
We instantiate (1) with a simple scalar problem  

\[
\min_{x\in\mathbb{R}} \; f(x) \;:=\;(x-3)^{2},
\qquad g\equiv0,\qquad K=1.
\]

Hence $g^{*}\equiv 0$, $\operatorname{prox}_{\sigma g^{*}}(u)=u$, and (3) reduces to  

\[
\begin{aligned}
y_{k+1}&=y_{k}+\sigma \bar{x}_{k},\\
x_{k+1}&=\operatorname{prox}_{\tau f}\!\bigl(x_{k}-\tau y_{k+1}\bigr).
\end{aligned}
\]

For $f$ quadratic,  

\[
\operatorname{prox}_{\tau f}(v)=\frac{v+2\tau\cdot3}{1+2\tau}
      =\frac{v+6\tau}{1+2\tau}.
\]

Choose $\tau=\sigma=0.1$, $\theta=1$, $x_{0}=0$, $y_{0}=0$, thus $\bar{x}_{0}=0$.

\begin{center}
\begin{tabular}{c|c|c|c|c}
$k$ & $y_{k+1}$ & $x_{k+1}$ & $\bar{x}_{k+1}$ & $f(x_{k+1})$\\\hline
0 & $0+0.1\cdot0=0$ &
$\displaystyle\frac{0-0.1\cdot0+0.6}{1.2}=0.5$ &
$0.5+1\cdot(0.5-0)=1.0$ &
$(0.5-3)^{2}=6.25$\\[2mm]
1 & $0+0.1\cdot1.0=0.1$ &
$\displaystyle\frac{0.5-0.1\cdot0.1+0.6}{1.2}
      =\frac{1.09}{1.2}=0.9083$ &
$0.9083+1\cdot(0.9083-0.5)=1.3166$ &
$(0.9083-3)^{2}=4.382$\\[2mm]
2 & $0.1+0.1\cdot1.3166=0.2317$ &
$\displaystyle\frac{0.9083-0.1\cdot0.2317+0.6}{1.2}
      =\frac{1.485}{1.2}=1.2375$ &
$1.2375+1\cdot(1.2375-0.9083)=1.5667$ &
$(1.2375-3)^{2}=3.095$\\
\end{tabular}
\end{center}

After three iterations $x_{3}\approx1.24$, already much closer to the optimum $x^{\star}=3$.

\section*{Assumptions}
\begin{assumption}[Convexity and Closedness]\label{ass:conv}
$f$ and $g$ are proper, closed, convex functions. Their conjugates $g^{*}$ are also convex and closed.
\end{assumption}
\begin{assumption}[Lipschitz Continuity of $K$]\label{ass:K}
$K$ is a bounded linear operator with operator norm $\|K\|<\infty$.
\end{assumption}
\begin{assumption}[Existence of a Saddle Point]\label{ass:saddle}
There exists $(x^{\star},y^{\star})$ such that  

\[
\mathcal{L}(x^{\star},y)\le\mathcal{L}(x^{\star},y^{\star})\le\mathcal{L}(x,y^{\star})\quad\forall x,y.
\]
\end{assumption}
\begin{assumption}[Step‑size Condition]\label{ass:steps}
The stepsizes satisfy  

\[
\tau\sigma\|K\|^{2}<1,\qquad \theta\in[0,1].
\]
\end{assumption}

\section*{Main Convergence Theorem}
\begin{theorem}[Ergodic $O(1/k)$ primal–dual gap]\label{thm:main}
Under Assumptions \ref{ass:conv}–\ref{ass:steps}, define the averaged iterates  

\[
\bar{x}^{N}:=\frac{1}{N}\sum_{k=0}^{N-1}x_{k+1},\qquad
\bar{y}^{N}:=\frac{1}{N}\sum_{k=0}^{N-1}y_{k+1}.
\]

Then for any $(x,y)$ we have  

\[
\mathcal{G}_{N}(x,y):=
\mathcal{L}(\bar{x}^{N},y)-\mathcal{L}(x,\bar{y}^{N})
\le\frac{1}{N}\Bigl(
\frac{\|x-x_{0}\|^{2}}{2\tau}
+\frac{\|y-y_{0}\|^{2}}{2\sigma}
\Bigr).
\tag{4}
\]

Consequently, choosing $(x,y)=(x^{\star},y^{\star})$ yields  

\[
\mathcal{G}_{N}(x^{\star},y^{\star})\le
\frac{C}{N},\qquad C:=\frac{\|x^{\star}-x_{0}\|^{2}}{2\tau}
            +\frac{\|y^{\star}-y_{0}\|^{2}}{2\sigma},
\]
i.e. an $O(1/N)$ convergence rate for the primal–dual gap.
\end{theorem}

\section*{Theoretical Properties}
\begin{itemize}
  \item \textbf{Stability.} Condition \eqref{ass:steps} guarantees that the Lyapunov function  
  $V_{k}:=\frac{1}{2\tau}\|x-x_{k}\|^{2}+\frac{1}{2\sigma}\|y-y_{k}\|^{2}$  
  is non‑increasing in expectation.
  \item \textbf{Hyperparameter Sensitivity.} The product $\tau\sigma$ must be strictly smaller than $1/\|K\|^{2}$; otherwise the descent inequality fails. Larger $\tau$ (resp. $\sigma$) speeds up primal (resp. dual) progress but tightens the admissible bound on the other stepsize.
  \item \textbf{Comparison.} For smooth $f$ and $g$, PDHG reduces to a forward–backward splitting scheme and recovers the $O(1/k)$ rate of the classical Chambolle–Pock method. In the purely smooth case it coincides with gradient descent with stepsize $\tau$ and thus inherits its optimal $O(1/k)$ rate for convex functions.
\end{itemize}

\section*{Discussion}
The theorem guarantees convergence of the *ergodic* (averaged) sequence. In practice, the raw iterates $(x_{k},y_{k})$ often converge faster. The extrapolation parameter $\theta=1$ yields the best known theoretical bound, but smaller $\theta$ can improve robustness for ill‑conditioned $K$.

\section*{Preliminaries (Definitions and Lemmas)}
\begin{lemma}[Three‑point identity for proximal operators]\label{lem:prox}
For any convex $h$ and $\tau>0$,
\[
\operatorname{prox}_{\tau h}(u)=\arg\min_{z}\Bigl\{h(z)+\frac{1}{2\tau}\|z-u\|^{2}\Bigr\}
\]
satisfies for all $z$,
\[
\langle z-\operatorname{prox}_{\tau h}(u),\, \operatorname{prox}_{\tau h}(u)-u\rangle
\ge\tau\bigl[h(\operatorname{prox}_{\tau h}(u))-h(z)\bigr].
\tag{5}
\]
\end{lemma}
\begin{proof}
Standard optimality condition of the proximal subproblem; see e.g. \cite{parikh2014prox}.
\end{proof}

\begin{lemma}[Descent lemma for the Lyapunov function]\label{lem:descent}
Let $(x_{k},y_{k})$ be generated by \eqref{3}. Then for any $(x,y)$,
\[
\begin{aligned}
&\frac{1}{2\tau}\|x-x_{k+1}\|^{2}
+\frac{1}{2\sigma}\|y-y_{k+1}\|^{2}\\
&\le
\frac{1}{2\tau}\|x-x_{k}\|^{2}
+\frac{1}{2\sigma}\|y-y_{k}\|^{2}
-\mathcal{L}(x_{k+1},y)+\mathcal{L}(x,y_{k+1})\\
&\quad+\frac{\theta}{2\tau}\|x_{k+1}-x_{k}\|^{2}
-\frac{(1-\tau\sigma\|K\|^{2})}{2\tau}\|x_{k+1}-\bar{x}_{k}\|^{2}.
\end{aligned}
\tag{6}
\]
\end{lemma}
\begin{proof}
We expand the squared norms using the update rules and apply Lemma \ref{lem:prox} twice (once for the dual, once for the primal).  
All algebraic steps are displayed in the main proof; the inequality \eqref{2} is used to drop the last non‑positive term.
\end{proof}

\section*{Main Proof (Step‑by‑Step Derivation)}
We prove Theorem \ref{thm:main}.  
Let $(x,y)$ be arbitrary. Define the Lyapunov function  

\[
V_{k}:=\frac{1}{2\tau}\|x-x_{k}\|^{2}+\frac{1}{2\sigma}\|y-y_{k}\|^{2}.
\tag{7}
\]

\noindent\textbf{[Step 1]} Apply Lemma \ref{lem:prox} to the dual update:
\[
\langle y-y_{k+1},\, y_{k+1}-(y_{k}+\sigma K\bar{x}_{k})\rangle
\ge\sigma\bigl[g^{*}(y_{k+1})-g^{*}(y)\bigr].
\tag{8}
\]

\noindent\textbf{[Step 2]} Rearrange \eqref{8} and use the definition of $\mathcal{L}$:
\[
g^{*}(y)-g^{*}(y_{k+1})\le
\frac{1}{\sigma}\langle y-y_{k+1},y_{k+1}-y_{k}\rangle
-\langle y-y_{k+1},K\bar{x}_{k}\rangle .
\tag{9}
\]

\noindent\textbf{[Step 3]} Apply Lemma \ref{lem:prox} to the primal update:
\[
\langle x-x_{k+1},\, x_{k+1}-(x_{k}-\tau K^{\!\top}y_{k+1})\rangle
\ge\tau\bigl[f(x_{k+1})-f(x)\bigr].
\tag{10}
\]

\noindent\textbf{[Step 4]} Rearrange \eqref{10}:
\[
f(x)-f(x_{k+1})\le
\frac{1}{\tau}\langle x-x_{k+1},x_{k+1}-x_{k}\rangle
+\langle x-x_{k+1},K^{\!\top}y_{k+1}\rangle .
\tag{11}
\]

\noindent\textbf{[Step 5]} Add \eqref{9} and \eqref{11} and collect the bilinear terms:
\[
\begin{aligned}
&\bigl[f(x)-f(x_{k+1})\bigr]+\bigl[g^{*}(y)-g^{*}(y_{k+1})\bigr]\\
&\le\frac{1}{\tau}\langle x-x_{k+1},x_{k+1}-x_{k}\rangle
      +\frac{1}{\sigma}\langle y-y_{k+1},y_{k+1}-y_{k}\rangle\\
&\quad+\langle x-x_{k+1},K^{\!\top}y_{k+1}\rangle
      -\langle y-y_{k+1},K\bar{x}_{k}\rangle .
\end{aligned}
\tag{12}
\]

\noindent\textbf{[Step 6]} Use the identity  
$\langle a,b\rangle=\frac{1}{2}\bigl(\|a\|^{2}+\|b\|^{2}-\|a-b\|^{2}\bigr)$
to rewrite the first two inner products:
\[
\begin{aligned}
\frac{1}{\tau}\langle x-x_{k+1},x_{k+1}-x_{k}\rangle
&=\frac{1}{2\tau}\bigl(\|x-x_{k}\|^{2}-\|x-x_{k+1}\|^{2}
                     -\|x_{k+1}-x_{k}\|^{2}\bigr),\\[1mm]
\frac{1}{\sigma}\langle y-y_{k+1},y_{k+1}-y_{k}\rangle
&=\frac{1}{2\sigma}\bigl(\|y-y_{k}\|^{2}-\|y-y_{k+1}\|^{2}
                     -\|y_{k+1}-y_{k}\|^{2}\bigr).
\end{aligned}
\tag{13}
\]

\noindent\textbf{[Step 7]} Expand the bilinear terms involving $K$:
\[
\begin{aligned}
\langle x-x_{k+1},K^{\!\top}y_{k+1}\rangle
   -\langle y-y_{k+1},K\bar{x}_{k}\rangle
&= \langle K(x-x_{k+1}),y_{k+1}\rangle
   -\langle K\bar{x}_{k},y-y_{k+1}\rangle\\
&= \langle K(x-\bar{x}_{k}),y_{k+1}\rangle
   -\langle K(x_{k+1}-\bar{x}_{k}),y_{k+1}\rangle
   -\langle K\bar{x}_{k},y-y_{k+1}\rangle\\
&= \langle K(x-\bar{x}_{k}),y_{k+1}\rangle
   -\langle K\bar{x}_{k},y-y_{k+1}\rangle
   -\langle K(x_{k+1}-\bar{x}_{k}),y_{k+1}\rangle .
\end{aligned}
\tag{14}
\]

\noindent\textbf{[Step 8]} Observe that  
$\langle K(x-\bar{x}_{k}),y_{k+1}\rangle-\langle K\bar{x}_{k},y-y_{k+1}\rangle
   =\langle Kx,y_{k+1}\rangle-\langle K\bar{x}_{k},y\rangle$.
Hence (12) becomes
\[
\begin{aligned}
&f(x)-f(x_{k+1})+g^{*}(y)-g^{*}(y_{k+1})\\
&\le
\frac{1}{2\tau}\bigl(\|x-x_{k}\|^{2}-\|x-x_{k+1}\|^{2}
                     -\|x_{k+1}-x_{k}\|^{2}\bigr)\\
&\quad+
\frac{1}{2\sigma}\bigl(\|y-y_{k}\|^{2}-\|y-y_{k+1}\|^{2}
                     -\|y_{k+1}-y_{k}\|^{2}\bigr)\\
&\quad+\langle Kx,y_{k+1}\rangle-\langle K\bar{x}_{k},y\rangle
 -\langle K(x_{k+1}-\bar{x}_{k}),y_{k+1}\rangle .
\end{aligned}
\tag{15}
\]

\noindent\textbf{[Step 9]} Add and subtract $\langle K\bar{x}_{k},y_{k+1}\rangle$ to the last line:
\[
-\langle K(x_{k+1}-\bar{x}_{k}),y_{k+1}\rangle
 =-\langle Kx_{k+1},y_{k+1}\rangle+\langle K\bar{x}_{k},y_{k+1}\rangle .
\tag{16}
\]

\noindent\textbf{[Step 10]} Combine the bilinear terms to obtain the primal–dual gap:
\[
\langle Kx,y_{k+1}\rangle-\langle K\bar{x}_{k},y\rangle
-\langle Kx_{k+1},y_{k+1}\rangle+\langle K\bar{x}_{k},y_{k+1}\rangle
= \mathcal{L}(x_{k+1},y)-\mathcal{L}(x,y_{k+1}) .
\tag{17}
\]

\noindent\textbf{[Step 11]} Insert \eqref{17} into \eqref{15} and rearrange:
\[
\begin{aligned}
\mathcal{L}(x_{k+1},y)-\mathcal{L}(x,y_{k+1})
&\le V_{k}-V_{k+1}
   -\frac{1}{2\tau}\|x_{k+1}-x_{k}\|^{2}
   -\frac{1}{2\sigma}\|y_{k+1}-y_{k}\|^{2}.
\end{aligned}
\tag{18}
\]

\noindent\textbf{[Step 12]} Use the extrapolation relation $\bar{x}_{k}=x_{k}+\theta(x_{k}-x_{k-1})$.
A standard algebraic manipulation (see Lemma 3.1 in \cite{chambolle2011first}) yields

\[
\|x_{k+1}-\bar{x}_{k}\|^{2}
= \|x_{k+1}-x_{k}\|^{2}
   -\theta(1-\theta)\|x_{k}-x_{k-1}\|^{2}.
\tag{19}
\]

Because $0\le\theta\le1$, the term $-\theta(1-\theta)\|x_{k}-x_{k-1}\|^{2}\le0$, and together with the stepsize condition \eqref{ass:steps} we obtain

\[
-\frac{1}{2\tau}\|x_{k+1}-x_{k}\|^{2}
\le
-\frac{1-\tau\sigma\|K\|^{2}}{2\tau}\|x_{k+1}-\bar{x}_{k}\|^{2}.
\tag{20}
\]

\noindent\textbf{[Step 13]} Plug \eqref{20} into \eqref{18} to get the descent inequality of Lemma \ref{lem:descent}:

\[
V_{k+1}\le V_{k}
-\bigl[\mathcal{L}(x_{k+1},y)-\mathcal{L}(x,y_{k+1})\bigr]
+\frac{\theta}{2\tau}\|x_{k+1}-x_{k}\|^{2}
-\frac{1-\tau\sigma\|K\|^{2}}{2\tau}\|x_{k+1}-\bar{x}_{k}\|^{2}.
\tag{21}
\]

Since the last two terms are non‑positive (by \eqref{ass:steps}), we finally have

\[
\mathcal{L}(x_{k+1},y)-\mathcal{L}(x,y_{k+1})
\le V_{k}-V_{k+1}.
\tag{22}
\]

\noindent\textbf{[Step 14]} Sum \eqref{22} for $k=0,\dots,N-1$:

\[
\sum_{k=0}^{N-1}\bigl[\mathcal{L}(x_{k+1},y)-\mathcal{L}(x,y_{k+1})\bigr]
\le V_{0}-V_{N}\le V_{0}.
\tag{23}
\]

\noindent\textbf{[Step 15]} Divide by $N$ and use convexity of $f$ and $g^{*}$ (which implies convexity of $\mathcal{L}$ in each argument) together with Jensen’s inequality:

\[
\mathcal{L}(\bar{x}^{N},y)-\mathcal{L}(x,\bar{y}^{N})
\le\frac{1}{N}\sum_{k=0}^{N-1}\bigl[\mathcal{L}(x_{k+1},y)-\mathcal{L}(x,y_{k+1})\bigr]
\le\frac{V_{0}}{N}.
\tag{24}
\]

Recalling the definition \eqref{7} of $V_{0}$ yields exactly the bound \eqref{4}, completing the proof of Theorem \ref{thm:main}. ∎

\section*{Rate Derivation (With Constants)}
From \eqref{4} and the definition of $V_{0}$ we have

\[
\mathcal{G}_{N}(x^{\star},y^{\star})
\le\frac{1}{N}
\left(
\frac{\|x^{\star}-x_{0}\|^{2}}{2\tau}
+\frac{\|y^{\star}-y_{0}\|^{2}}{2\sigma}
\right)
=: \frac{C}{N}.
\]

Thus the primal–dual gap decays with rate $O(1/N)$. The constant $C$ depends linearly on the squared distances of the initial point to the optimal saddle point and inversely on the stepsizes.

\section*{Special Cases}
\begin{itemize}
  \item \textbf{Purely primal problem ($g\equiv0$).} Then $y_{k}\equiv0$, and \eqref{3} reduces to a proximal gradient step
    $x_{k+1}=\operatorname{prox}_{\tau f}(x_{k})$, recovering the classical $O(1/k)$ rate for convex nonsmooth minimization.
  \item \textbf{Both $f$ and $g$ smooth.} The proximal operators become explicit gradient steps, and the algorithm coincides with the Arrow–Hurwicz method with over‑relaxation; the same $O(1/k)$ bound holds.
  \item \textbf{Choice $\theta=0$.} No extrapolation; inequality \eqref{19} simplifies and the convergence proof still goes through, albeit with a slightly larger constant $C$.
\end{itemize}

\bibliographystyle{plain}
\begin{thebibliography}{9}
\bibitem{parikh2014prox}
N. Parikh and S. Boyd, \emph{Proximal Algorithms}, Foundations and Trends in Optimization, 2014.
\bibitem{chambolle2011first}
A. Chambolle and T. Pock, \emph{A First-Order Primal-Dual Algorithm for Convex Problems with Applications to Imaging}, J. Math. Imaging Vis., 2011.
\end{thebibliography}

\end{document}